\documentclass[fontsize=15pt]{jlreq}

%プリアンブル
\usepackage{amsmath}
\pagestyle{empty}
\usepackage[top=15mm, bottom=20mm, left=20mm, right=20mm]{geometry}

%本文
\begin{document}
\begin{center}
{\LARGE {振り返り}} \\
-ダニエル電池- %単元ごとに変えるか消す
\end{center}
\vspace{1em}
\begin{flushright}
名前：\underline{\hspace{5cm}}
\end{flushright}

% 問題部分
\begin{enumerate}
    \item 電圧計は回路に対して直列に繋ぐか，並列に繋ぐか．\\
    ( \hspace{2cm} )
    \vspace{2em} 

    \item ダニエル電池では，+極に何の金属を使用するか．\\
    ( \hspace{2cm} )
    \vspace{2em} 

    \item ダニエル電池ではなぜセロハンが必要か．\\
    ( \hspace{10cm} )
    \vspace{2em} 

    \item 電池とは，物質が持つ(\hspace{2cm})エネルギーを(\hspace{2cm})エネルギーに変換する装置．
    \vspace{2em} 

    \item 一次電池，二次電池とはそれぞれどのような電池か．
    \begin{itemize}
      \item 一次電池：
      \vspace{2em}
      \item 二次電池：
    \end{itemize}
    \vspace{2em} 

    \item 一次電池，二次電池の例をそれぞれひとつずつ挙げよ．
    \begin{itemize}
      \item 一次電池：
      \vspace{2em}
      \item 二次電池：
    \end{itemize}

\end{enumerate}

\end{document}
